\section{Threats to Validity}
\textbf{Internal validity.} S1 shares advisory inputs, model setting, nonblank source eligibility, evaluation, and acquisition ceiling. Policy-specific prompts, interaction, actual acquired information, and cumulative input remain different. Native completion behavior and empty-answer handling are part of Agent's configuration. Association between exhaustion and failure is not causal evidence. S0 additionally has eligibility/index-cap asymmetries; S1 bundles multiple changes after inspecting S0 and reuses its sample. The study estimates configuration differences and does not isolate the effect of an acquisition policy or repair.

\textbf{Construct validity.} The task uses dataset-supplied entry titles alongside public descriptions. Entry titles can contain annotation-added filenames; 47 titles match at least one target basename. Explicit label-field withholding does not remove this semantic information. All configurations share it, so it changes task difficulty and interpretation rather than candidate-universe fairness. The description-only flags cannot identify hint-free inputs. Basename matching need not be unique, and its association with scores would not estimate removal effects. Public-code training overlap is unmeasured. FR@5 measures unique target-file coverage and ignores false positives. LR@5 rewards tolerant interval overlap, while LR@5$_{\leq64}$ only limits interval width. Neither validates a complete EP--CO trace, explanation, exploit, or discovery of an unknown defect. The prompt does not separately request EP and CO, so role recall measures annotation coverage rather than full task compliance. Schema validation checks integer coordinates but not file existence or bounds; the additional location checks are descriptive. Trajectory-wide source unions do not identify when or how the model used evidence. Human-verified benchmark labels were not independently reannotated in this study. Pack's weak lexical prefix limits the meaning of comparisons against it.

\textbf{Statistical validity.} The 80 entries have distinct advisories but are clustered in 13 repositories with substantial imbalance: 45 entries come from Openclaw. Entry-level sign flips and bootstrap intervals assume independent differences and may understate uncertainty if within-repository outcomes are correlated. The post hoc repository sign-flip sensitivity also needs independent, sign-exchangeable repository contributions. Its separately adjusted secondary probabilities illustrate dependence sensitivity without certifying a general population effect. The 13 uneven clusters support limited-precision diagnostics; leave-out and repository-resampling summaries do not establish cluster-robust validity. Official secondary probabilities follow the frozen four-test BH rule; intervals are unadjusted. S0 and S1 are not pooled, and a nonsignificant contrast is not interpreted as equivalence.

\textbf{External validity.} Findings concern one hosted model, one information ceiling, a selected verified subset with dataset-title assistance, and three specific implementations. Transfer to original-advisory inputs without these title additions is unknown. Three repeats describe observed variability without controlling the provider's random generation or hidden backend. A more capable Pack, iterative RAG, other agent scaffolds, independent advisories, and additional models may differ. S1's known sample limits its use as external confirmation.

\textbf{Measurement and provenance.} The source ledger excludes Pi argument-validation feedback and provider history replay. $A$ and $B$ use different accounting scopes/tokenizers. Online time excludes initial validation and some loading, while local indexing has no complete amortized cost estimate. Reported charges depend on usage/cache accounting; two historical charges remain unresolved. The saved Pi events expose final model behavior but not complete historical HTTP bytes or an immutable archive of the whole SDK. Attribution between provider and adapter therefore remains unresolved. Source/record/hash consistency establishes an inspectable evidence chain within these boundaries, not absence of every implementation defect.
